\chapter{In Which Aramis is Sought, and Only  Bazin is Found}

\lettrine{T}{wo} hours had scarcely elapsed since the departure of  the master of the house, who, in Blaisois's sight, had taken the road to  Paris, when a horseman, mounted on a good pied horse, stopped before the gate,  and with a sonorous <hola!> called the stable-boys, who, with the  gardeners, had formed a circle round Blaisois, the historian-in-ordinary to the  household of the château. This <hola,> doubtless well known to  Master Blaisois, made him turn his head and exclaim—<Monsieur  d'Artagnan! run quickly, you chaps, and open the gate.>

A swarm of eight brisk lads flew to the gate, which was opened as if it had  been made of feathers; and every one loaded him with attentions, for they knew  the welcome this friend was accustomed to receive from their master; and for  such remarks the eye of the valet may always be depended upon.

<Ah!> said M. d'Artagnan, with an agreeable smile, balancing  himself upon his stirrup to jump to the ground, <where is that dear  count?>

<Ah! how unfortunate you are, monsieur!> said Blaisois: <and  how unfortunate will monsieur le comte, our master, think himself when he hears  of your coming! As ill luck will have it, monsieur le comte left home two hours  ago.>

D'Artagnan did not trouble himself about such trifles. <Very  good!> said he. <You always speak the best French in the world; you  shall give me a lesson in grammar and correct language, whilst I wait the  return of your master.>

<That is impossible, monsieur,> said Blaisois; <you would  have to wait too long.>

<Will he not come back to-day, then?>

<No, nor to-morrow, nor the day after to-morrow. Monsieur le comte has  gone on a journey.>

<A journey!> said D'Artagnan, surprised; <that's  a fable, Master Blaisois.>

<Monsieur, it is no more than the truth. Monsieur has done me the honour  to give me the house in charge; and he added, with his voice so full of  authority and kindness—that is all one to me: <You will say I have  gone to Paris.>>

<Well!> cried D'Artagnan, <since he is gone towards  Paris, that is all I wanted to know! you should have told me so at first,  booby! He is then two hours in advance?>

<Yes, monsieur.>

<I shall soon overtake him. Is he alone?>

<No, monsieur.>

<Who is with him, then?>

<A gentleman whom I don't know, an old man, and M. Grimaud.>

<Such a party cannot travel as fast as I can—I will start.>

<Will monsieur listen to me an instant?> said Blaisois, laying his  hand gently on the reins of the horse.

<Yes, if you don't favour me with fine speeches, and make  haste.>

<Well, then, monsieur, that word Paris appears to me to be only an  excuse.>

<Oh, oh!> said D'Artagnan, seriously, <an excuse,  eh?>

<Yes, monsieur: and monsieur le comte is not going to Paris, I will  swear.>

<What makes you think so?>

<This,—M. Grimaud always knows where our master is going; and he  had promised me that the first time he went to Paris, he would take a little  money for me to my wife.>

<What, have you a wife, then?>

<I had one—she was of this country; but monsieur thought her a  noisy scold, and I sent her to Paris; it is sometimes inconvenient, but very  agreeable at others.>

<I understand; but go on. You do not believe the count gone to  Paris?>

<No, monsieur; for then M. Grimaud would have broken his word; he would  have perjured himself, and that is impossible.>

<That is impossible,> repeated D'Artagnan, quite in a study,  because he was quite convinced. <Well, my brave Blaisois, many thanks to  you.>

Blaisois bowed.

<Come, you know I am not curious—I have serious business with your  master. Could you not, by a little bit of a word—you who speak so  well—give me to understand—one syllable only—I will guess the  rest.>

<Upon my word, monsieur, I cannot. I am quite ignorant where monsieur le  comte is gone. As to listening at doors, that is contrary to my nature; and  besides, it is forbidden here.>

<My dear fellow,> said D'Artagnan, <this is a very bad  beginning for me. Never mind; you know when monsieur le comte will return, at  least?>

<As little, monsieur, as the place of his destination.>

<Come, Blaisois, come, search.>

<Monsieur doubts my sincerity? Ah, monsieur, that grieves me much.>

<The devil take his gilded tongue!> grumbled D'Artagnan.  <A clown with a word would be worth a dozen of him. Adieu!>

<Monsieur, I have the honour to present you my respects.>

<Cuistre!> said D'Artagnan to himself, <the fellow is  unbearable.> He gave another look up to the house, turned his  horse's head, and set off like a man who has nothing either annoying or  embarrassing in his mind. When he was at the end of the wall, and out of  sight,—<Well, now, I wonder,> said he, breathing quickly,  <whether Athos was at home. No; all those idlers, standing with their  arms crossed, would have been at work if the eye of the master was near. Athos  gone on a journey?—that is incomprehensible. Bah! it is all devilish  mysterious! And then—no—he is not the man I want. I want one of a  cunning, patient mind. My business is at Melun, in a certain presbytery I am  acquainted with. Forty-five leagues—four days and a half! Well, it is  fine weather, and I am free. Never mind the distance!>

And he put his horse into a trot, directing his course towards Paris. On the  fourth day he alighted at Melun, as he had intended.

D'Artagnan was never in the habit of asking any one on the road for any  common information. For these sorts of details, unless in very serious  circumstances, he confided in his perspicacity, which was so seldom at fault,  in his experience of thirty years, and in a great habit of reading the  physiognomies of houses, as well as those of men. At Melun, D'Artagnan  immediately found the presbytery—a charming house, plastered over red  brick, with vines climbing along the gutters, and a cross, in carved stone,  surmounting the ridge of the roof. From the ground-floor of this house came a  noise, or rather a confusion of voices, like the chirping of young birds when  the brood is just hatched under the down. One of these voices was spelling the  alphabet distinctly. A voice thick, yet pleasant, at the same time scolded the  talkers and corrected the faults of the reader. D'Artagnan recognized  that voice, and as the window of the ground-floor was open, he leant down from  his horse under the branches and red fibres of the vine and cried,  <Bazin, my dear Bazin! good-day to you.>

A short, fat man, with a flat face, a cranium ornamented with a crown of gray  hairs, cut short, in imitation of a tonsure, and covered with an old black  velvet cap, arose as soon as he heard D'Artagnan—we ought not to  say arose, but bounded up. In fact, Bazin bounded up, carrying with him his  little low chair, which the children tried to take away, with battles more  fierce than those of the Greeks endeavouring to recover the body of Patroclus  from the hands of the Trojans. Bazin did more than bound; he let fall both his  alphabet and his ferule. <You!> said he; <you, Monsieur  D'Artagnan?>

<Yes, myself! Where is Aramis—no, M. le Chevalier  d'Herblay—no, I am still mistaken—Monsieur le  Vicaire-General?>

<Ah, monsieur,> said Bazin, with dignity, <monseigneur is at  his diocese.>

<What did you say?> said D'Artagnan. Bazin repeated the  sentence.

<Ah, ah! but has Aramis a diocese?>

<Yes, monsieur. Why not?>

<Is he a bishop, then?>

<Why, where can you come from,> said Bazin, rather irreverently,  <that you don't know that?>

<My dear Bazin, we pagans, we men of the sword, know very well when a man  is made a colonel, or maitre-de-camp, or marshal of France; but if he be made a  bishop, arch-bishop, or pope—devil take me if the news reaches us before  the three quarters of the earth have had the advantage of it!>

<Hush! hush!> said Bazin, opening his eyes: <do not spoil  these poor children, in whom I am endeavouring to inculcate such good  principles.> In fact, the children had surrounded D'Artagnan, whose  horse, long sword, spurs, and martial air they very much admired. But above  all, they admired his strong voice; so that, when he uttered his oath, the  whole school cried out, <The devil take me!> with fearful bursts of  laughter, shouts, and bounds, which delighted the musketeer, and bewildered the  old pedagogue.

<There!> said he, <hold your tongues, you brats! You have  come, M. d'Artagnan, and all my good principles fly away. With you, as  usual, comes disorder. Babel is revived. Ah! Good Lord! Ah! the wild little  wretches!> And the worthy Bazin distributed right and left blows which  increased the cries of his scholars by changing the nature of them.

<At least,> said he, <you will no longer decoy any one  here.>

<Do you think so?> said D'Artagnan, with a smile which made a  shudder creep over the shoulders of Bazin.

<He is capable of it,> murmured he.

<Where is your master's diocese?>

<Monseigneur Rene is bishop of Vannes.>

<Who had him nominated?>

<Why, monsieur le surintendant, our neighbour.>

<What! Monsieur Fouquet?>

<To be sure he did.>

<Is Aramis on good terms with him, then?>

<Monseigneur preached every Sunday at the house of monsieur le  surintendant at Vaux; then they hunted together.>

<Ah!>

<And monseigneur composed his homilies—no, I mean his  sermons—with monsieur le surintendant.>

<Bah! he preached in verse, then, this worthy bishop?>

<Monsieur, for the love of heaven, do not jest with sacred things.>

<There, Bazin, there! So, then, Aramis is at Vannes?>

<At Vannes, in Bretagne.>

<You are a deceitful old hunks, Bazin; that is not true.>

<See, monsieur, if you please; the apartments of the presbytery are  empty.>

<He is right there,> said D'Artagnan, looking attentively at  the house, the aspect of which announced solitude.

<But monseigneur must have written you an account of his  promotion.>

<When did it take place?>

<A month back.>

<Oh! then there is no time lost. Aramis cannot yet have wanted me. But  how is it, Bazin, you do not follow your master?>

<Monsieur, I cannot; I have occupations.>

<Your alphabet?>

<And my penitents.>

<What, do you confess, then? Are you a priest?>

<The same as one. I have such a call.>

<But the orders?>

<Oh,> said Bazin, without hesitation, <now that monseigneur  is a bishop, I shall soon have my orders, or at least my dispensations.>  And he rubbed his hands.

<Decidedly,> said D'Artagnan to himself, <there will be  no means of uprooting these people. Get me some supper, Bazin.>

<With pleasure, monsieur.>

<A fowl, a bouillon, and a bottle of wine.>

<This is Saturday night, monsieur—it is a day of abstinence.>

<I have a dispensation,> said D'Artagnan.

Bazin looked at him suspiciously.

<Ah, ah, master hypocrite!> said the musketeer, <for whom do  you take me? If you, who are the valet, hope for dispensation to commit a  crime, shall not I, the friend of your bishop, have dispensation for eating  meat at the call of my stomach? Make yourself agreeable with me, Bazin, or by  heavens! I will complain to the king, and you shall never confess. Now you know  that the nomination of bishops rests with the king,—I have the king, I am  the stronger.>

Bazin smiled hypocritically. <Ah, but we have monsieur le  surintendant,> said he.

<And you laugh at the king, then?>

Bazin made no reply; his smile was sufficiently eloquent.

<My supper,> said D'Artagnan, <it is getting towards  seven o'clock.>

Bazin turned round and ordered the eldest of the pupils to inform the cook. In  the meantime, D'Artagnan surveyed the presbytery.

<Phew!> said he, disdainfully, <monseigneur lodged his  grandeur very meanly here.>

<We have the château de Vaux,> said Bazin.

<Which is perhaps equal to the Louvre?> said D'Artagnan,  jeeringly.

<Which is better,> replied Bazin, with the greatest coolness  imaginable.

<Ah, ah!> said D'Artagnan.

He would perhaps have prolonged the discussion, and maintained the superiority  of the Louvre, but the lieutenant perceived that his horse remained fastened to  the bars of a gate.

<The devil!> said he. <Get my horse looked after; your master  the bishop has none like him in his stables.>

Bazin cast a sidelong glance at the horse, and replied, <Monsieur le  surintendant gave him four from his own stables; and each of the four is worth  four of yours.>

The blood mounted to the face of D'Artagnan. His hand itched and his eye  glanced over the head of Bazin, to select the place upon which he should  discharge his anger. But it passed away; reflection came, and D'Artagnan  contented himself with saying,—

<The devil! the devil! I have done well to quit the service of the king.  Tell me, worthy Master Bazin,> added he, <how many musketeers does  monsieur le surintendant retain in his service?>

<He could have all there are in the kingdom with his money,>  replied Bazin, closing his book, and dismissing the boys with some kindly blows  of his cane.

<The devil! the devil!> repeated D'Artagnan, once more, as if  to annoy the pedagogue. But as supper was now announced, he followed the cook,  who introduced him into the refectory, where it awaited him. D'Artagnan  placed himself at the table, and began a hearty attack upon his fowl.

<It appears to me,> said D'Artagnan, biting with all his  might at the tough fowl they had served up to him, and which they had evidently  forgotten to fatten,—<it appears that I have done wrong in not  seeking service with that master yonder. A powerful noble this intendant,  seemingly! In good truth, we poor fellows know nothing at the court, and the  rays of the sun prevent our seeing the large stars, which are also suns, at a  little greater distance from our earth,—that is all.>

As D'Artagnan delighted, both from pleasure and system, in making people  talk about things which interested him, he fenced in his best style with Master  Bazin, but it was pure loss of time; beyond the tiresome and hyperbolical  praises of monsieur le surintendant of the finances, Bazin, who, on his side,  was on his guard, afforded nothing but platitudes to the curiosity of  D'Artagnan, so that our musketeer, in a tolerably bad humour, desired to  go to bed as soon as he had supped. D'Artagnan was introduced by Bazin  into a mean chamber, in which there was a poor bed; but D'Artagnan was  not fastidious in that respect. He had been told that Aramis had taken away the  key of his own private apartment, and as he knew Aramis was a very particular  man, and had generally many things to conceal in his apartment, he had not been  surprised. He, therefore, although it seemed comparatively even harder,  attacked the bed as bravely as he had done the fowl; and, as he had as good an  inclination to sleep as he had had to eat, he took scarcely longer time to be  snoring harmoniously than he had employed in picking the last bones of the  bird.

Since he was no longer in the service of any one, D'Artagnan had promised  himself to indulge in sleeping as soundly as he had formerly slept lightly; but  with whatever good faith D'Artagnan had made himself this promise, and  whatever desire he might have to keep it religiously, he was awakened in the  middle of the night by a loud noise of carriages, and servants on horseback. A  sudden illumination flashed over the walls of his chamber; he jumped out of bed  and ran to the window in his shirt. <Can the king be coming this  way?> he thought, rubbing his eyes; <in truth, such a suite can  only be attached to royalty.>

<Vive le monsieur le surintendant!> cried, or rather vociferated,  from a window on the ground-floor, a voice which he recognized as  Bazin's, who at the same time waved a handkerchief with one hand, and  held a large candle in the other. D'Artagnan then saw something like a  brilliant human form leaning out of the principal carriage; at the same time  loud bursts of laughter, caused, no doubt, by the strange figure of Bazin, and  issuing from the same carriage, left, as it were, a train of joy upon the  passage of the rapid cortège.

<I might easily see it was not the king,> said D'Artagnan;  <people don't laugh so heartily when the king passes. Hola,  Bazin!> cried he to his neighbour, three-quarters of whose body still hung  out of the window, to follow the carriage with his eyes as long as he could.  <What is all that about?>

<It is M. Fouquet,> said Bazin, in a patronizing tone.

<And all those people?>

<That is the court of M. Fouquet.>

<Oh, oh!> said D'Artagnan; <what would M. de Mazarin  say to that if he heard it?> And he returned to his bed, asking himself  how Aramis always contrived to be protected by the most powerful personages in  the kingdom. <Is it that he has more luck than I, or that I am a greater  fool than he? Bah!> That was the concluding word by the aid of which  D'Artagnan, having become wise, now terminated every thought and every  period of his style. Formerly he said, <Mordioux!> which was a  prick of the spur, but now he had become older, and he murmured that  philosophical <Bah!> which served as a bridle to all the passions.  